\documentclass[a4paper,11pt]{article}
\usepackage{booktabs}
\usepackage{array}
\usepackage{xcolor}
\usepackage{colortbl}
\usepackage{amsmath}
\usepackage{geometry}
\usepackage{caption}
\geometry{margin=25mm}

\definecolor{best}{rgb}{0.85,0.95,0.85}
\definecolor{zero}{rgb}{0.95,0.93,0.93}
\definecolor{partial}{rgb}{0.98,0.98,0.88}

\title{Comparison of Three Selection Rules on Archimedean Solids\\
       \large Success Rates of the Spiral Unfolding Algorithm R1/R2/R3}
\author{}
\date{April 4, 2026}

\begin{document}
\maketitle

\section{Overview}

We applied three face-selection rules of the spiral unfolding algorithm to all 13 Archimedean solids
and compared the unfolding success rates over all pairs $(C_1, C_2)$.
The orientation was fixed to Face-up (the topmost face centroid aligned with the $+z$ direction).

\bigskip
\noindent
Definition of each rule:
\begin{description}
  \item[R1 (max $\varphi$)] Among the left candidates, select the face with the largest angle $\varphi$ from the forward direction (sharpest left turn). If no left candidate exists, select the face with the largest $|\varphi|$.
  \item[R2 (loxodrome)] Among the left candidates, select the face with the smallest $|\varphi|$ (closest to straight ahead). This approximates a loxodrome on the sphere. If no left candidate exists, select the face with the smallest $|\varphi|$.
  \item[R3 (max $z$)] Among the left candidates, select the face whose centroid has the highest $z$-coordinate. If no left candidate exists, select the face with the lowest $z$.
\end{description}

\section{Results}

\begin{table}[h]
\centering
\caption{Unfolding success rates by rule for the 13 Archimedean solids (Face-up orientation)}
\label{tab:archimedean}
\renewcommand{\arraystretch}{1.25}
\begin{tabular}{lrrrrrrrrr}
\toprule
Polyhedron & Faces & Pairs
  & R1 & R1\% & R2 & R2\% & R3 & R3\% \\
\midrule
\rowcolor{partial}
Truncated Tetrahedron       &  8 &  36 & 12  & 33.3 &  5 & 13.9 & 11 & 30.6 \\
\rowcolor{zero}
Cuboctahedron               & 14 &  48 &  0  &  0.0 &  0 &  0.0 &  0 &  0.0 \\
\rowcolor{zero}
Truncated Cube              & 14 &  72 &  0  &  0.0 &  0 &  0.0 &  0 &  0.0 \\
\rowcolor{best}
Truncated Octahedron        & 14 &  72 & \textbf{64}  & \textbf{88.9} &  0 &  0.0 & 23 & 31.9 \\
\rowcolor{zero}
Rhombicuboctahedron         & 26 &  96 &  0  &  0.0 &  0 &  0.0 &  0 &  0.0 \\
\rowcolor{zero}
Truncated Cuboctahedron     & 26 & 144 &  0  &  0.0 &  3 &  2.1 &  1 &  0.7 \\
\rowcolor{partial}
Snub Cube                   & 38 & 120 & 18  & 15.0 & \textbf{47} & \textbf{39.2} & 20 & 16.7 \\
\rowcolor{zero}
Icosidodecahedron           & 32 & 120 &  0  &  0.0 &  0 &  0.0 &  0 &  0.0 \\
\rowcolor{zero}
Truncated Dodecahedron      & 32 & 180 &  0  &  0.0 &  0 &  0.0 &  0 &  0.0 \\
\rowcolor{best}
Truncated Icosahedron       & 32 & 180 & \textbf{154} & \textbf{85.6} &  0 &  0.0 & 26 & 14.4 \\
\rowcolor{zero}
Rhombicosidodecahedron      & 62 & 240 &  0  &  0.0 &  0 &  0.0 &  0 &  0.0 \\
\rowcolor{zero}
Truncated Icosidodecahedron & 62 & 360 &  0  &  0.0 &  0 &  0.0 &  0 &  0.0 \\
\rowcolor{zero}
Snub Dodecahedron           & 92 & 300 &  0  &  0.0 &  0 &  0.0 &  0 &  0.0 \\
\bottomrule
\end{tabular}
\end{table}

\noindent
{\footnotesize Row colors: \colorbox{best}{green} = high success rate (best rule $\geq 80\%$),
\colorbox{partial}{yellow} = partial success, \colorbox{zero}{red} = 0\% for all rules.}

\section{Discussion}

\subsection{Majority of Solids Yield 0\% for All Rules}

Nine of the 13 solids (69\%) achieved a success rate of 0\% for all rules.
These are concentrated among polyhedra containing many square faces
(Cuboctahedron, rhombic families, Truncated Cuboctahedron, Snub Dodecahedron, etc.).
The mixture of face types is thought to disrupt the continuity of the spiral unfolding.

\subsection{Solids Where R1 Is Effective}

R1 substantially outperformed the other rules for the Truncated Tetrahedron (33.3\%),
Truncated Octahedron (88.9\%), and Truncated Icosahedron (85.6\%).
These polyhedra consist exclusively of regular triangles, hexagons, or pentagons,
and their homogeneous face structure appears well-suited to spiral unfolding.

Notably, R1 achieved 100\% on the regular Dodecahedron (Platonic),
whereas it dropped to 0\% on the Truncated Dodecahedron (faces: regular pentagons and decagons).
The large decagonal faces introduced by truncation are conjectured to exhaust the pool of
faces satisfying the left-turn condition.

\subsection{Reversal in the Snub Cube}

The Snub Cube was the only solid where R2 (39.2\%) $>$ R1 (15.0\%).
The asymmetric arrangement of squares and triangles in snub polyhedra may incidentally
favor loxodrome-style selection (R2).
However, the Snub Dodecahedron yielded 0\% for all rules,
suggesting that this reversal may be specific to snub polyhedra with a smaller face count.

\subsection{Generally Low Success Rate of R2}

R2 yielded 0\% for the Truncated Octahedron and Truncated Icosahedron,
the very solids where R1 performs best.
This is consistent with the behavior on the regular Dodecahedron
(where R2 always halts at 6 faces---the ``6-face wall'' phenomenon),
suggesting that the loxodrome selection rule triggers a similar stalling problem
even in structures dominated by hexagons and pentagons.

\section{Comparison with Platonic Solids}

\begin{table}[h]
\centering
\caption{Unfolding success rates by rule for the Platonic solids (reference)}
\label{tab:platonic}
\renewcommand{\arraystretch}{1.25}
\begin{tabular}{lrrrr}
\toprule
Polyhedron & Faces & R1\% & R2\% & R3\% \\
\midrule
Tetrahedron   &  4 & 100.0 & 100.0 & 100.0 \\
Cube          &  6 & 100.0 & 100.0 & 100.0 \\
Octahedron    &  8 & 100.0 & 100.0 & 100.0 \\
Dodecahedron  & 12 & 100.0 &   0.0 &  47--53 \\
Icosahedron   & 20 & 100.0 & 100.0 & 100.0 \\
\bottomrule
\end{tabular}
\end{table}

While R1 achieves 100\% on all five Platonic solids,
it is effective for only 3 of the 13 Archimedean solids and peaks at 88.9\%.
This indicates that uniformity of face types plays a decisive role
in the success of spiral unfolding.

\end{document}
